The problem of realisability of global models, e.g., global (session)
types \cite{BEJLERI-2009,Castagna-2011,Honda-2016,Hans-2016}, global
choreographies \cite{Franco-2022}, global labelled transition
systems \cite{Beek-2023} and global languages
has attracted many researchers over the past several years.
A realisation of a global model is a set of
collaborating computing entities which
are executed concurrently, and whose behaviours, obtained from their
(synchronous or asynchronous) interactions, are exactly those that are
obtained from the global models.
For example, a global model can be the parallel composition of a set
of labelled transition systems, each of which models the local
behaviour of a component.
An interesting, but challenging question is how to determine whether a
global model is realisable and, if so, how to obtain a realisation.

We are interested in the problem of realisability for global
models that take the form of sets of sequential timed
scenarios. Given such a set, our goal is to
obtain a \emph{timed scenario expression} that realises the set.

In earlier work we developed the notions of sequential
timed scenarios (scenarios for short) \cite{neda-2018-2} and
distributed timed scenarios \cite{DTS-2025}.
A scenario is a sequence of (names of) events along with a set of
constraints on the times at which these events occur
(see \Sec{sec-seq} for more details).
A distributed timed scenario (DTS) is a collection of
components, each of which is a timed
scenario (see \Sec{sec-dts} for more details).
Each component performs its specified actions, represented as a
sequence of events, concurrently with
other components. Components can communicate or
synchronise with each other by performing the same action
simultaneously.
%The idea is that each component performs its specified
%actions, represented as a sequence of events, independently from
%other components. However, from time to
%time, the components communicate or synchronise with each other
%by performing the same action simultaneously.
A DTS, $\DScenario$, is viewed as a parallel composition of its components and its
semantics is defined as a set of distributed behaviours, i.e., the set
of all possible interleavings of events in the components of
$\DScenario$ that satisfy the constraints and the synchronising
requirements of $\DScenario$.

Given a set, $\mathcal{S}$, of scenarios over a set of events,
$\Sigma_{\mathcal{S}}$,
the stable distance relation \cite{DTS-2025} of $\mathcal{S}$
is $\DR{\mathcal{S}}=(\ParS{\mathcal{S}},\DF{\mathcal{S}})$.
The partial order $\ParS{\mathcal{S}}$
captures all the ordering dependencies between events in
$\mathcal{S}$, while the distance function, $\DF{\mathcal{S}}$,
expresses constraints on the time between
every pair of events $e$ and $e'$,
where $e\ParS{\mathcal{S}}e'$. A projection of
$\DR{\mathcal{S}}$ is a scenario over $\Sigma_{\mathcal{S}}$ whose
sequence of events agrees with
$\ParS{\mathcal{S}}$ and whose constraints satisfy
%includes the tightest constraints that satisfy
the constraints captured by $\DF{\mathcal{S}}$.

In the current paper we introduce \emph{timed scenario expressions}
and solve the problem of realisability of sets of
scenarios as scenario expressions.
%That is, given a global model of the behaviours of a system, captured
%by a set of scenarios, $\mathcal{S}$, we produce a scenario expression that expresses
%the same set of behaviours that are allowed by the members of $\mathcal{S}$.
So our contribution is twofold:
%\subsubsection{Contributions}
% \setlist{nolistsep}
%\begin{enumerate}[noitemsep]
%\item

First, we introduce a limited form of timed scenario expressions (expressions
for short), as a means
of constructing hierarchical models of complex systems, whose constituent
parts are distributed systems (\Sec{sec-exp}). Intuitively, an expression is built by
repeatedly applying three operators, namely choice,
concatenation and parallel composition (denoted by $\BAR$, $\circ$ and
$||$\footnote{We do not consider the Kleene star ($*$) operator here, as its
introduction opens up interesting questions that cannot
be addressed in the confines of the current paper.}, respectively), to
distributed scenarios. In its simplest form an
expression is just a single DTS.
Given an expression $E$, we
define its semantics, denoted by $\Sem{E}$, in terms of the set
of distributed behaviours that are expressed by $E$.
%\item

Second, we tackle the problem of realisability of sets of (sequential)
scenarios as expressions (\Sec{sec-realize}). Given a finite
set, $\mathcal{S}$, of scenarios, we produce a scenario
expression, $E$, that realises $\mathcal{S}$,
i.e., $\Sem{E}$ is identical to the union of the sets of behaviours
specified by the members of $\mathcal{S}$.

We identify two criteria for realisability of $\mathcal{S}$ as a
single DTS: $\mathcal{S}$ is realisable by a DTS iff (i) it
is \emph{order-closed}, and (ii) all its members are \emph{compatible}, in
terms of their constraints, with $\DR{\mathcal{S}}$.
Intuitively,
$\mathcal{S}$ is order-closed if, for every projection of
$\DR{\mathcal{S}}$, there is a scenario in $\mathcal{S}$
with that exact sequence of events, and vice versa.
A scenario $\SeqSC\in\mathcal{S}$ is compatible with
$\DR{\mathcal{S}}$ if its set of constraints is equivalent to the set
of constraints of the projection of $\DR{\mathcal{S}}$ that
has the same sequence of events as $\SeqSC$.

If $\mathcal{S}$ is not realisable by a single DTS,
then it must be divided into
partitions, in such a way that each partition is order-closed and all
its members are compatible with the distance relation of the
partition. Once every partition is realisable by a single DTS, the
resulting expression will be a choice between these
distributed scenarios.

The novelty of our work is primarily in the use of (stable) distance
relations \cite{DTS-2025} as a linchpin of an efficient method
%, confirmed by our experimental results,
for the realisation problem.
By identifying the necessary and sufficient conditions for the
realisability of a set of scenarios as a DTS, we address the problem
that was left open in our earlier work \cite{DTS-2025}.



%The reason we do not consider the Kleene star ($*$) operator in our
%expressions in the current paper is that they make the realisation
%problem significantly more complex.
%To fully address the problem of realisability where the
%resulting expressions involve $*$
%would require developing new apparatus and
%techniques (in addition to what is developed in the current paper),
%which could not be presented in the current paper due to
%the page limitation.
%\end{enumerate}
%\vspace{-1.4em}
\subsubsection{Related Work}
In our earlier work \cite{DTS-2025} we developed an algorithm for
solving the problem
of realisability of sets of (sequential) scenarios as distributed scenarios.
Given a set of scenarios, our algorithm determines whether
the set is realisable by a \emph{single} DTS, i.e.,
whether there is a DTS whose semantics is identical
to the union of the behaviours described by the members of
the set. If so, the algorithm produces such a DTS.
Despite some similarities, the problem that we tackle in the current
paper is fundamentally different.

%----------------------
The comparison of our sequential timed scenarios
with simple temporal networks (STNs) \cite{DECHTER1991}, as well as
with message sequence charts (MSCs) \cite{Harel2003,MSCs}, is presented
elsewhere \cite{DTS-2025}.
Similarities between
%our distance relations and timed partial
%orders (TPOs) of Watanabe et al. \cite{watanabe2023}, as well as
%between
the distributed behaviours of distributed timed scenarios and the
timed traces of Balaguer et al. \cite{hal-2012} are also discussed
in our earlier work \cite{DTS-2025}.

%The comparison of our stable distance tables with Difference Bounds
%Matrices is presented elsewhere \cite{neda-2017,neda-2018-2}.

The problem of realisability has been investigated in a variety of settings,
%of global models
%has been investigated
e.g., for message sequence
charts \cite{Alur-2003}, global session
types \cite{BEJLERI-2009,Castagna-2011,Honda-2016,Hans-2016},
pomsets \cite{GUANCIALE201969},
global labelled transition systems \cite{Beek-2023},
%more recently for
choreography languages \cite{Franco-2022} and team
automata \cite{Beek-2024}. Unlike in all
these works the realisability of our sets of sequential scenarios as
scenario expressions is
determined by precise time analysis enabled by our distance
relations.
%The distance relations enable us to reason about the underlying partail
%order, as well as the time constraints between various events in
%scenarios, both of which are essential for realisability of
That is, in our setting we are able to directly reason
about realisability of models that include explicit time constraints
and produce realisations that involve time constraints.

